\documentclass[a4paper,11pt]{article}
\usepackage[utf8]{inputenc}
\usepackage{geometry}
\usepackage{xcolor}
\usepackage{listings}
\usepackage{hyperref}
\usepackage{graphicx}
\usepackage{amsmath}
\usepackage{amssymb}
\usepackage{tcolorbox}
\usepackage{float}
\usepackage{booktabs}

% Page geometry
\geometry{top=2.5cm, bottom=2.5cm, left=2.5cm, right=2.5cm}

% Custom colours
\definecolor{ibmblue}{RGB}{0, 45, 156}
\definecolor{ibmred}{RGB}{250, 75, 76}
\definecolor{purpleaccent}{RGB}{109, 40, 217}
\definecolor{amberaccent}{RGB}{245, 158, 11}
\definecolor{tealaccent}{RGB}{13, 148, 136}
\definecolor{shockorange}{RGB}{234, 88, 12}
\definecolor{codegreen}{rgb}{0,0.6,0}
\definecolor{codegray}{rgb}{0.5,0.5,0.5}
\definecolor{codepurple}{rgb}{0.58,0,0.82}
\definecolor{backcolour}{rgb}{0.95,0.95,0.92}

% Python code style
\lstdefinestyle{mystyle}{
    backgroundcolor=\color{backcolour},
    commentstyle=\color{codegreen},
    keywordstyle=\color{magenta},
    numberstyle=\tiny\color{codegray},
    stringstyle=\color{codepurple},
    basicstyle=\ttfamily\footnotesize,
    breakatwhitespace=false,
    breaklines=true,
    captionpos=b,
    keepspaces=true,
    numbers=left,
    numbersep=5pt,
    showspaces=false,
    showstringspaces=false,
    showtabs=false,
    tabsize=2
}
\lstset{style=mystyle}

% Custom boxes (titles always in braced form)
\newtcolorbox{studentbox}{colback=blue!5!white,colframe=blue!75!black,title={Student Activity}}
\newtcolorbox{theorybox}{colback=gray!5!white,colframe=gray!50!black,title={Theory Recap}}
\newtcolorbox{warningbox}{colback=yellow!10!white,colframe=orange!80!black,title={Important}}
\newtcolorbox{simulatorbox}{colback=purple!5!white,colframe=purple!75!black,title={Simulator}}
\newtcolorbox{reflectionbox}{colback=teal!5!white,colframe=teal!75!black,title={Quick Check}}
\newtcolorbox{connectionbox}{colback=green!5!white,colframe=green!60!black,title={Connections to Previous Labs}}
\newtcolorbox{procedurebox}{colback=red!4!white,colframe=red!70!black,title={Procedure 1 vs Procedure 2}}
\newtcolorbox{questionbox}{colback=violet!4!white,colframe=violet!70!black,title={The Question}}
\newtcolorbox{shockbox}{colback=shockorange!8!white,colframe=shockorange,title={Counterintuitive --- and Real}}
\newtcolorbox{benchbox}{colback=teal!6!white,colframe=tealaccent,title={Open the Simulator First}}
\newtcolorbox{numberbox}{colback=green!4!white,colframe=green!50!black,title={Formula $\rightarrow$ Number $\rightarrow$ Bench}}
\newtcolorbox{industrybox}{colback=amberaccent!10!white,colframe=amberaccent!90!black,title={Real-World Anchor}}
\newtcolorbox{deepdivebox}{colback=cyan!4!white,colframe=cyan!60!black,title={Deep Dive}}

% Bra-ket shorthands (guarded: never clash with package-provided macros)
\providecommand{\ket}[1]{|#1\rangle}
\providecommand{\bra}[1]{\langle#1|}
\providecommand{\braket}[2]{\langle#1|#2\rangle}
\providecommand{\ketbra}[2]{|#1\rangle\langle#2|}

\title{\textbf{Laboratory Guide 7: Quantum Optics --- Photons, Interference, and Entanglement}\\
\large The Hardware of Quantum Technology}
\author{Course: Quantum Computing Foundation}
\date{}

\begin{document}

\maketitle

\section{Introduction and Objectives}

Labs 5 and 6 gave you the mathematics of qubits and a protocol built on it. This lab
shows you the \textbf{hardware those abstractions run on}. A single photon \emph{is} a
qubit --- twice over: its \textbf{path} through an interferometer is one two-level
system, its \textbf{polarization} is another. Every device on the optical bench ---
beam splitter, phase shifter, polarizer --- is one of the operators you already know,
executed by glass and crystal. And every strange number this lab produces (a
coincidence counter pinned at zero, two 50/50 elements composing to certainty, a
filter that \emph{revives} blocked light, twin photons that always agree) is a
Born-rule computation you have already done abstractly.

\vspace{0.3cm}
\noindent\textbf{Prerequisites:}
\begin{itemize}
    \item \textit{Fundamental Concepts of Quantum Technologies} (lecture series).
    \item Lab 5: \textit{Quantum Measurements} --- the Born rule, measurement
          Procedures 1 and 2, and the states $\ket{0}, \ket{1}, \ket{+}, \ket{-}$
          on the Bloch sphere.
    \item Lab 6: \textit{Quantum Key Distribution (BB84)} --- mutual unbiasedness
          and the measurement--disturbance--security bridge.
    \item Familiarity with Dirac notation: $\ket{\psi}$, $\bra{\phi}$, $\braket{\phi}{\psi}$.
\end{itemize}

\begin{connectionbox}
\textbf{The Lab 5--6 bridge.} The path states $\ket{A}, \ket{B}$ and the polarization
states $\ket{H}, \ket{V}$ are both the Lab 5 computational ($Z$) basis; the diagonal
polarizations at $\pm 45^\circ$ are exactly the Lab 5 $X$-basis states
$\ket{+}, \ket{-}$. The four BB84 signal states of Lab 6 are \emph{physically}
$\ket{H}, \ket{V}, \ket{+}, \ket{-}$ --- Part C of this lab is the transmitter and
receiver hardware of a real QKD link, and Part B's which-path result is Lab 6's
``Eve leaves fingerprints'' theorem restated as an identity about interference.
\end{connectionbox}

\noindent\textbf{Mode:} Interactive Optical Bench + Qiskit \quad\textbf{Credits:} 1 ECTS
\quad\textbf{Estimated time:} $\approx$ 10 h

\vspace{0.3cm}
\noindent\textbf{Learning Objectives:}
\begin{enumerate}
    \item \textbf{Explain the operational principles of single-photon and
          entangled-photon experiments} --- beam splitters as unitaries, detection as
          projective measurement, anti-correlation, interference, and coincidence
          counting.
    \item \textbf{Design and analyze simple optical setups} --- predict the output
          statistics of a Mach--Zehnder interferometer, a polarizer cascade, or a
          two-analyzer coincidence experiment, and verify the prediction photon by
          photon in the simulator.
    \item \textbf{Interpret interference and coincidence data} --- fringe visibility,
          which-path complementarity ($\mathcal{V} = s$), and the quantitative gap
          between quantum and classical-wave predictions.
    \item \textbf{Discuss the role of optical platforms in quantum communication and
          computation} --- polarization-encoded and entanglement-based QKD,
          interferometric sensing, and dual-rail photonic quantum computing.
\end{enumerate}

\begin{simulatorbox}
The lab's interactive page (\texttt{lab7\_optics.html}) contains the
\textbf{Optical Bench} built for this lab: five guided presets reproduce exactly the
experiments of this guide --- \textbf{Preset A} (beam splitter + coincidence counter),
\textbf{Preset B1} (Mach--Zehnder fringes), \textbf{Preset B2} (which-path marker),
\textbf{Preset C} (polarizer cascade), \textbf{Preset D} (entangled pairs, optional).
Every displayed probability comes from the same Born-rule pipeline as the formulas in
this guide, and all runs are seeded: the same seed and the same sequence of actions
always reproduce the same counts.

\medskip
\textbf{This guide is written simulator-first.} Each Part opens with an ``Open the
Simulator First'' box. Do that two-minute run \emph{before} reading the theory: the
counts you observe pose exactly the question the section then answers. The companion
Qiskit notebook rebuilds the same optics as quantum circuits:\\[4pt]
\url{https://colab.research.google.com/github/LMrilo/quantum-labs-2026/blob/main/lab7/lab7_optics.ipynb}
\end{simulatorbox}

\newpage

\section{Part A: One Photon at a Beam Splitter}

\begin{benchbox}
Open the bench in \textbf{Preset A}: a single-photon source, one 50/50 beam splitter,
a detector on each output port, and a coincidence counter wired to both. Fire
\textbf{10\,000 photons} and watch two things: the two detector tallies, and the
coincidence counter. The tallies drift toward 50/50 --- but the coincidence counter
\emph{never moves}. Keep that second observation in mind; it is the whole section.
\end{benchbox}

\begin{questionbox}
What does it actually mean that light comes in photons --- and what happens when one
indivisible photon hits a half-silvered mirror?
\end{questionbox}

\begin{shockbox}
One photon \textbf{takes both paths} --- the state after the beam splitter is a
superposition over the two output arms, and Part B will prove it by interference ---
yet it \textbf{never triggers both detectors}. Not rarely: \emph{never}. The photon
is divisible in amplitude but indivisible in detection.
\end{shockbox}

A half-silvered mirror (a \textbf{50/50 beam splitter}) transmits half the light and
reflects half. For a classical wave that sentence is unambiguous: the energy splits,
half into each arm. But light arrives in indivisible quanta. A single photon cannot
deliver half a click to each detector --- so what does ``splits 50/50'' mean for
\emph{one photon}?

Quantum mechanics answers with a state, not a story. Write the two output paths as
the basis states $\ket{A}$ (transmitted) and $\ket{B}$ (reflected) --- the same
two-level system as any Lab 5 qubit. A lossless beam splitter must act as a
\textbf{unitary} on this qubit (unitarity here \emph{is} energy conservation --- the
photon must exit somewhere), and the symmetric 50/50 choice is:

\begin{equation}
    B = \frac{1}{\sqrt{2}}\begin{pmatrix} 1 & i \\ i & 1 \end{pmatrix}
    \qquad\Longrightarrow\qquad
    B\ket{A} = \frac{\ket{A} + i\,\ket{B}}{\sqrt{2}}.
    \label{eq:bs}
\end{equation}

\begin{theorybox}
\textbf{Why the $i$?} A naive ``50/50 without phase'',
$\tfrac{1}{\sqrt 2}\bigl(\begin{smallmatrix}1&1\\1&1\end{smallmatrix}\bigr)$, is a
singular matrix --- it is \emph{not a possible physical device}, because it would
merge two distinguishable inputs into one output and lose a photon in the process.
The phase $i$ on reflection is the minimal repair that restores unitarity
(Zeilinger, 1981). This phase convention --- $i$ on reflection, used consistently ---
is fixed for the entire lab: every formula below is derived under it, and the
simulator uses exactly this matrix.
\end{theorybox}

Detection at the two output ports is a path-basis measurement ---
\textbf{Procedure 1} from Lab 5, executed destructively (the photon is absorbed by
the detector that clicks). The Born rule gives the two outcome probabilities:

\begin{equation}
    P(\text{detector at port A}) = \left|\tfrac{1}{\sqrt 2}\right|^2 = \tfrac12,
    \qquad
    P(\text{detector at port B}) = \left|\tfrac{i}{\sqrt 2}\right|^2 = \tfrac12,
    \label{eq:bsprobs}
\end{equation}
\begin{equation}
    P(\text{both detectors}) = 0 \quad\text{exactly.}
    \label{eq:anticorr}
\end{equation}

That last line deserves a hard look. The state $(\ket{A} + i\ket{B})/\sqrt 2$ is
\textbf{one photon in superposition over two paths --- not two photons}. A path
measurement has exactly two mutually exclusive outcomes; ``both detectors click'' is
not a low-probability event, it is \emph{not in the outcome set at all}. The zero is
structural, not statistical.

\begin{numberbox}
$P = \tfrac12, \tfrac12$ and $P_{\text{coinc}} = 0$. In \textbf{Preset A}: after
10\,000 photons the tallies land within statistical noise of
\textbf{5\,000 / 5\,000} (the theory overlay shows the exact 50\% line), while the
coincidence counter reads \textbf{0 --- exactly, after every single photon}. Fire
another 10\,000 and check that the zero is not luck.
\end{numberbox}

\begin{warningbox}
\textbf{The honest classical contrast.} A classical wave or pulse really does split:
half the energy in each arm, and with two detectors both can fire in the same trial.
The standard figure of merit is $g^{(2)}(0)$, the normalized zero-delay second-order
correlation --- the coincidence rate normalized to random chance (Loudon, 2000).
Semiclassical models predict coincidence rates at least at the random-chance level,
$g^{(2)}(0) \geq 1$, while the one-photon state gives $g^{(2)}(0) = 0$. Grangier,
Roger and Aspect measured exactly this anti-correlation in 1986; it is the
operational meaning of ``the photon is indivisible'', and it is the first of this
lab's four experimental pillars.
\end{warningbox}

\begin{industrybox}
\textbf{Where this is an industry, not a thought experiment.} One photon distributed
over two spatial modes is the \textbf{dual-rail qubit} --- the native encoding of
photonic quantum computing. \textbf{Quandela} builds the single-photon sources (the
``fire one photon on demand'' box the simulator idealizes) and dual-rail photonic
processors around exactly this $\{\ket{A}, \ket{B}\}$ qubit; \textbf{ORCA Computing}
builds photonic quantum processors on the same principle that beam-splitter networks
perform unitary computation on light. The linear-optical quantum computing
literature (the KLM protocol) showed that such networks plus detection are universal.
When you fire photons at the bench, you are operating the elementary cell of those
machines.
\end{industrybox}

\section{Part B: The Mach--Zehnder Interferometer and Which-Path Information}

\begin{benchbox}
Switch the bench to \textbf{Preset B1}: two beam splitters back to back, mirrors
folding the two arms together, and a phase knob $\varphi$ on arm B. Leave
$\varphi = 0^\circ$ and fire 10\,000 photons. Every photon exits at \textbf{D1} ---
100/0, from a device built of two 50/50 elements. Now sweep the $\varphi$ slider and
watch the split move. The counts \emph{depend on the phase difference between the two
paths}.
\end{benchbox}

\begin{questionbox}
If the photon went one way or the other, the detector counts could not depend on the
path it didn't take. They do. What now?
\end{questionbox}

\begin{shockbox}
Two 50/50 beam splitters in sequence give a \textbf{100/0} outcome. Model each
splitter as an independent coin flip --- the reasonable move if the photon ``really''
takes one path --- and you predict 50/50 at the output. The observed perfect nulls
and perfect peaks are the single-photon interference signature: the amplitudes of
\emph{both} paths add before anything is detected.
\end{shockbox}

The \textbf{Mach--Zehnder interferometer} is the workhorse: beam splitter BS1 opens
two paths, mirrors bring them back together, a \textbf{phase shifter}
$\Phi(\varphi) = \mathrm{diag}(1, e^{i\varphi})$ on arm B models any path-length
difference, and beam splitter BS2 recombines. (The mirrors add the same phase to both
arms --- a global phase, dropped.) Step the input $\ket{A}$ through the device:

\begin{equation}
    \ket{A}
    \;\xrightarrow{\;\text{BS1}\;}\; \frac{\ket{A} + i\ket{B}}{\sqrt 2}
    \;\xrightarrow{\;\Phi(\varphi)\;}\; \frac{\ket{A} + i e^{i\varphi}\ket{B}}{\sqrt 2}
    \;\xrightarrow{\;\text{BS2}\;}\;
    \frac{(1 - e^{i\varphi})\ket{A} + i(1 + e^{i\varphi})\ket{B}}{2}.
    \label{eq:mzpipe}
\end{equation}

With detector \textbf{D1} on output port B and \textbf{D2} on output port A, the Born
rule and the identity $|1 \pm e^{i\varphi}|^2 = 2 \pm 2\cos\varphi$ give the fringes:

\begin{equation}
    P(D1) = \cos^2\!\frac{\varphi}{2}, \qquad P(D2) = \sin^2\!\frac{\varphi}{2}.
    \label{eq:fringes}
\end{equation}

At $\varphi = 0$: $P(D1) = 1$ --- the 100/0 you just observed. Algebraically this is
the composition identity $B^2 = iX$: two symmetric beam splitters are a perfect path
swap (up to an overall phase). At $\varphi = \pi$ the same instrument, retuned, sends
every photon to D2 instead. Between the extremes the split interpolates as a perfect
cosine:

\begin{table}[H]
\centering
\begin{tabular}{lcccc}
\toprule
phase $\varphi$ & $0^\circ$ & $60^\circ$ & $90^\circ$ & $180^\circ$ \\
\midrule
$P(D1) = \cos^2(\varphi/2)$ & $1$ & $3/4$ & $1/2$ & $0$ \\
$P(D2) = \sin^2(\varphi/2)$ & $0$ & $1/4$ & $1/2$ & $1$ \\
\bottomrule
\end{tabular}
\caption{Mach--Zehnder anchor points, equation~\eqref{eq:fringes}. The
$60^\circ$ column is the 3\,:\,1 fringe point used by Preset B1.}
\label{tab:mz}
\end{table}

For fringes of the form $P(\varphi) = (1 + \mathcal{V}\cos\varphi)/2$, the
\textbf{fringe visibility} is
\begin{equation}
    \mathcal{V} = \frac{P_{\max} - P_{\min}}{P_{\max} + P_{\min}},
    \label{eq:vis}
\end{equation}
and the unmarked interferometer has $P_{\max} = 1$, $P_{\min} = 0$:
$\mathcal{V} = 1$.

\begin{numberbox}
At $\varphi = \pi/3$ ($60^\circ$): $P(D1) = \cos^2(30^\circ) = \tfrac34$ --- a
\textbf{75\% / 25\%} split. In \textbf{Preset B1}, set the slider to $60^\circ$ and
fire 10\,000 photons: the tallies land on 75/25 within noise, on top of the theory
overlay. Then return to $\varphi = 0^\circ$ and re-verify the \textbf{100/0} --- two
50/50 splitters, certain outcome.
\end{numberbox}

\begin{industrybox}
\textbf{The fringe as a measuring instrument.} $\varphi$ is optical path-length
difference in disguise. Any physical quantity that stretches one arm --- rotation in
a \textbf{fiber-optic gyroscope} (the navigation-grade rotation sensors in aircraft),
strain, temperature, or a passing gravitational wave in the kilometer-scale
\textbf{LIGO-class interferometers} --- shifts the fringe. Near the steepest point of
the cosine, reading $P(D1)$ to precision $\delta P$ reads the phase to
$\delta\varphi \approx 2\,\delta P$: counting photons \emph{is} the measurement. That
is interferometric sensing, and it is this exact circuit.
\end{industrybox}

\subsection{Which-Path Marking: Complementarity, Quantified}

\begin{benchbox}
Switch to \textbf{Preset B2}: the same interferometer, plus a \textbf{which-path
marker} in the arms --- a probe qubit that records which arm the photon took. Its
dial is the \textbf{marker overlap $s$}: at $s = 1$ the marker states for the two
arms are identical (no record); at $s = 0$ they are orthogonal (a perfect record).
Set $s = 0$, sweep $\varphi$, and watch the fringes: they are \textbf{gone} --- flat
50/50 at every phase. Now slide $s$ upward and watch them come back, exactly in
proportion.
\end{benchbox}

\begin{shockbox}
An \textbf{unread} which-path marker still destroys the fringes. Nobody looks at the
record. It is enough that the record \emph{exists} --- that the universe could in
principle tell which path was taken. Interference is lost the moment the paths become
distinguishable, not the moment someone learns the answer.
\end{shockbox}

Model the marking honestly. Between the beam splitters, couple the path qubit to a
marker qubit: path A tags the marker into state $\ket{M_A}$, path B into $\ket{M_B}$
(normalized, with overlap $s = \braket{M_A}{M_B}$ taken real in $[0,1]$). We write
the two-qubit state as $\ket{\text{path}\;\text{marker}}$ with the \textbf{path qubit
leftmost} --- big-endian ordering, qubit 0 is the leftmost symbol; the companion
notebook converts Qiskit's opposite bitstring order before displaying anything.
After BS1 and the marking:

\begin{equation}
    \frac{\ket{A}\ket{M_A} + i e^{i\varphi}\,\ket{B}\ket{M_B}}{\sqrt 2}.
    \label{eq:marked}
\end{equation}

Recombine at BS2 (acting on the path qubit only) and compute $P(D1)$, summing over
the marker --- read or unread, the arithmetic cannot tell:

\begin{equation}
    P(D1; \varphi, s)
    = \frac{1}{4}\,\bigl\|\, \ket{M_A} + e^{i\varphi}\ket{M_B} \,\bigr\|^2
    = \frac{1 + s\cos\varphi}{2}
    \qquad\Longrightarrow\qquad
    \boxed{\;\mathcal{V} = s = |\braket{M_A}{M_B}|\;}
    \label{eq:complementarity}
\end{equation}

\textbf{Fringe visibility equals the indistinguishability of the marker states.}
$s = 1$: no record, full fringes. $s = 0$: perfect record, no fringes ---
$P(D1) = \tfrac12$ at every phase, as if the photon really had flipped a coin at BS1.
In between, the trade is exactly linear in the fringe term. A perfect marking
($s = 0$) is physically a controlled-NOT --- the path controls the marker, turning
the superposition into the entangled path--marker state
$(\ket{00} + i e^{i\varphi}\ket{11})/\sqrt 2$ (path leftmost, as always) --- and
reading the marker afterwards is a genuine path measurement
(\textbf{Procedure 1} on the marker qubit): same statistics as putting a detector in
the arm, but the photon survives to reach BS2.

\begin{numberbox}
At $s = 0.5$: $\mathcal{V} = 0.5$, so the fringes run between
$P_{\max} = \tfrac{1+0.5}{2} = \mathbf{0.75}$ and
$P_{\min} = \tfrac{1-0.5}{2} = \mathbf{0.25}$. In \textbf{Preset B2}, set $s = 0.5$,
fire at $\varphi = 0^\circ$ ($\rightarrow$ 75\%) and at $\varphi = 180^\circ$
($\rightarrow$ 25\%), and read the visibility card: it converges to 0.5 --- the dial
you set, recovered from counting statistics alone.
\end{numberbox}

\begin{connectionbox}
Lab 5 said ``measurement disturbs the state''; Lab 6 said ``Eve's measurement leaves
fingerprints''. Both are this one identity $\mathcal{V} = |\braket{M_A}{M_B}|$: an
eavesdropper who gains path information ($s < 1$) \emph{necessarily} degrades the
interference any honest party could observe. Information out $\Longrightarrow$
coherence lost, in exact proportion. (The general form is Englert's duality relation
$\mathcal{V}^2 + \mathcal{D}^2 \leq 1$ with distinguishability
$\mathcal{D} = \sqrt{1 - s^2}$ --- our ideal marker saturates it.)
\end{connectionbox}

\section{Part C: Polarization --- Quantum Malus and the Three-Polarizer Surprise}

\begin{benchbox}
Switch to \textbf{Preset C}: an unpolarized source and a rack of polarizing filters.
Start with two \textbf{crossed} polarizers --- $0^\circ$ and $90^\circ$. Fire
10\,000 photons: \textbf{nothing gets through}. Now insert a $45^\circ$ polarizer
between them and fire again: \textbf{12.5\%} of the photons pass all three filters.
You added an absorber, and the light came back.
\end{benchbox}

\begin{questionbox}
Two crossed filters block everything. Insert a third filter between them and light
comes through. Why does adding an absorber \emph{increase} transmission?
\end{questionbox}

\begin{shockbox}
Adding an absorbing polarizer raises transmission from exactly \textbf{0} to
\textbf{1/8}. No sieve, filter, or ``the polarizer lets through the right
components'' picture survives this: a filter that only ever \emph{removes} light
cannot make more light arrive. Something other than removal is happening. (You can
try this at home with three lenses from polarized sunglasses.)
\end{shockbox}

Polarization is the photon's second qubit: $\ket{H}$ (horizontal) and $\ket{V}$
(vertical) form the computational basis, and the state at linear angle $\theta$ from
the H-axis is $\ket{\theta} = \cos\theta\,\ket{H} + \sin\theta\,\ket{V}$. The
diagonal states at $\pm 45^\circ$ are exactly the Lab 5 $X$-basis pair $\ket{\pm}$.

An ideal polarizer at angle $\theta$ is the projector
$\Pi_\theta = \ketbra{\theta}{\theta}$, and this is the crucial fact: a polarizer is
an inherently \textbf{Procedure-1 device}. Each photon either passes --- \emph{and is
then in state $\ket{\theta}$} --- or is absorbed. There is no fractional photon. For
a photon in $\ket{H}$ meeting a polarizer at $\theta$, the Born rule gives the
\textbf{quantum Malus law}:

\begin{equation}
    P(\text{pass}) = |\braket{\theta}{H}|^2 = \cos^2\theta,
    \qquad \text{post-state } \ket{\theta}.
    \label{eq:malus}
\end{equation}

\begin{table}[H]
\centering
\begin{tabular}{lccccc}
\toprule
polarizer angle $\theta$ & $0^\circ$ & $30^\circ$ & $45^\circ$ & $60^\circ$ & $90^\circ$ \\
\midrule
$P(\text{pass})$ from $\ket{H}$ & $1$ & $3/4$ & $1/2$ & $1/4$ & $0$ \\
\bottomrule
\end{tabular}
\caption{Quantum Malus law, equation~\eqref{eq:malus}, at the five anchor angles.}
\label{tab:malus}
\end{table}

Classical Malus ($I = I_0\cos^2\theta$) reads identically --- but the quantum version
is \textbf{per-photon and probabilistic}: $\cos^2\theta$ is a Born probability, the
same $|\braket{m}{\psi}|^2$ as every measurement since Lab 5, and each individual
photon entirely passes or is entirely absorbed.

\begin{numberbox}
$\cos^2(60^\circ) = \mathbf{25\%}$. In \textbf{Preset C}, switch the source to
$\ket{H}$, enable a single polarizer at $60^\circ$, and fire 10\,000 photons: the
pass fraction lands on 25\% within noise. Then rebuild the crossed pair (unpolarized
source, $0^\circ$ and $90^\circ$) $\rightarrow$ \textbf{exactly 0} --- and insert the
$45^\circ$ filter $\rightarrow$ \textbf{1/8 = 12.5\%}.
\end{numberbox}

Now the resolution of the three-polarizer surprise. For a cascade
$\theta_1, \dots, \theta_N$, each passage \emph{renews} the state to the current
polarizer's $\ket{\theta_j}$ --- the Procedure-1 post-state. So the pass
probabilities multiply with the \emph{differences} between consecutive angles:

\begin{equation}
    P(\text{pass all}) = P_1 \cdot \prod_{j=2}^{N}\cos^2(\theta_j - \theta_{j-1}),
    \label{eq:cascade}
\end{equation}

\noindent with $P_1 = \tfrac12$ for an unpolarized source (no polarization
preference --- every first filter passes half). The two experiments you ran:

\begin{equation}
    \text{crossed } (0^\circ, 90^\circ):\;\;
    \tfrac12\cdot\cos^2(90^\circ) = 0,
    \qquad\quad
    \text{three } (0^\circ, 45^\circ, 90^\circ):\;\;
    \tfrac12\cdot\tfrac12\cdot\tfrac12 = \tfrac18.
    \label{eq:threepol}
\end{equation}

The middle filter does not ``let more light through'' --- it \textbf{rewrites the
state}. After the $45^\circ$ polarizer the photon is in $\ket{+}$, which is no longer
orthogonal to the final $90^\circ$ filter. The projective post-state --- the fact
that measurement \emph{changes} the system, the disturbance you met in Lab 5 --- is
the entire explanation. The $45^\circ$ polarizer is an $X$-basis measurement of the
polarization qubit, executed by a plastic film.

\begin{procedurebox}
\textbf{Two physically different devices.} Lab 5 distinguished two ways to measure in
the diagonal basis, and in optics they are \emph{different hardware}:
\begin{itemize}
    \item \textbf{Procedure 1} --- a polarizer at $45^\circ$: probabilities
          $|\braket{\pm}{\psi}|^2$, passed photon left in $\ket{+}$ (the $\ket{-}$
          outcome is an absorbed photon).
    \item \textbf{Procedure 2} --- a half-waveplate at $22.5^\circ$ (which implements
          the canonical basis-change rotation $U_X = H$) followed by a polarizing
          beam splitter and two detectors: \emph{identical probabilities}, but the
          photon exits in $\ket{H}$ or $\ket{V}$, not $\ket{\pm}$.
\end{itemize}
Same statistics, different post-measurement states: the overlap between the two
devices' ``+'' outputs is $|\braket{+}{H}|^2 = \tfrac12$, not 1. Downstream optics
see different photons. Lab 6's naive-resend attack --- where Eve's forgotten
re-encoding turned a 25\% error rate into 37.5\% --- was exactly this distinction
carrying a security consequence.
\end{procedurebox}

\begin{industrybox}
\textbf{This hardware runs QKD networks today.} Polarization-encoded BB84 uses
precisely these components: Lab 6's four signal states \emph{are}
$\ket{H}, \ket{V}, \ket{+}, \ket{-}$; Alice's transmitter is a polarization
preparer, Bob's receiver is a basis-switched analyzer --- a rotation plus a
polarizing beam splitter, Procedure 2 in glass. Metropolitan QKD fiber networks
(deployed in Europe and Asia, and the building block of the planned
\textbf{EuroQCI} infrastructure) run this Part's optics at telecom scale: every
``secure key bit'' started life as a photon meeting a polarization analyzer.
\end{industrybox}

\begin{deepdivebox}
\textbf{The polarizer staircase.} Generalize the insertion trick: between $0^\circ$
and $90^\circ$, insert intermediate filters in equal steps of $90^\circ/N$ --- $N$
steps, hence $N+1$ polarizers counting both ends. The transmission
$\tfrac12\cos^{2N}(90^\circ/N)$ \emph{rises} toward $\tfrac12$ as $N$ grows: many
gentle projections approximate a lossless rotation of the polarization. ($N = 2$ is
the three-polarizer experiment.) Try the staircase
$0^\circ, 30^\circ, 60^\circ, 90^\circ$ in Preset C:
$\tfrac12\,(3/4)^3 = 27/128 \approx 21.1\%$, already far above the 12.5\% of the
single insertion.
\end{deepdivebox}

\section{Part D: Entangled Photon Pairs \textnormal{\normalsize(Optional)}}

\noindent\textit{This part is optional enrichment: nothing in the activities' core
depends on it. It is also the doorway to the deepest experimental result in the
field --- read on if you want to see it.}

\begin{benchbox}
Switch to \textbf{Preset D}: a crystal that emits \textbf{pairs} of photons, one to
each side, each side ending in an analyzer (angle dial + two detectors, ``+'' and
``$-$''). Set both dials to the \textbf{same angle} --- any angle --- and fire
10\,000 pairs. Each side alone clicks ``+'' or ``$-$'' like a fair coin. But the
\textbf{agreement} counter reads \textbf{100\%}: the two sides always give the same
answer. Now change the common angle and fire again --- still 100\%.
\end{benchbox}

\begin{questionbox}
Two photons, each individually a perfect coin toss, always agree when asked the same
question --- at any angle, chosen after they separated. What kind of correlation is
that?
\end{questionbox}

\begin{shockbox}
Entangled photons are \textbf{individually random at every angle, perfectly
correlated whenever the angles match} --- and the best classical wave model is not
just vaguely wrong about this, it is off by \textbf{exactly a factor of 2} in the
correlation amplitude. The simulator plots both curves; the gap between them is the
measurement.
\end{shockbox}

Nonlinear crystals (spontaneous parametric down-conversion, the workhorse
entangled-pair source since the 1990s --- Kwiat et al., 1995) emit two photons in the
joint polarization state

\begin{equation}
    \ket{\Phi^{+}} = \frac{\ket{HH} + \ket{VV}}{\sqrt 2},
    \label{eq:pairstate}
\end{equation}

\noindent written photon 1 leftmost, photon 2 rightmost --- the same leftmost-first,
big-endian ordering as Part B's $\ket{\text{path}\;\text{marker}}$ kets. Each side
passes a two-channel analyzer at angle $\alpha$ (photon 1) or $\beta$ (photon 2):
outcome ``+'' means the photon exited the $\ket{\theta_{\text{analyzer}}}$ port,
``$-$'' the orthogonal port. Two facts, both exact:

\begin{itemize}
    \item \textbf{Each photon alone: a perfect coin.} Summing over the partner's
          outcomes leaves photon 1 (or 2) in the maximally mixed state $I/2$ ---
          \emph{exactly} unpolarized. $P(+) = \tfrac12$ at \emph{every} analyzer
          angle. No analyzer setting on one side is visible in the other side's
          single rates.
    \item \textbf{The pair together: locked.} The joint amplitude is
          $\braket{\alpha, \beta}{\Phi^{+}}
          = (\cos\alpha\cos\beta + \sin\alpha\sin\beta)/\sqrt 2
          = \cos(\alpha - \beta)/\sqrt 2$, so with $\Delta = \alpha - \beta$:
          \begin{equation}
              P({+}{+}) = P({-}{-}) = \frac{\cos^2\Delta}{2},
              \qquad
              P({+}{-}) = P({-}{+}) = \frac{\sin^2\Delta}{2}.
              \label{eq:jointprobs}
          \end{equation}
          At matched angles ($\Delta = 0$): agreement with certainty. At
          $\alpha = 0^\circ$, $\beta = 22.5^\circ$, for instance:
          $P({+}{+}) = P({-}{-}) = (1 + \cos 45^\circ)/4 \approx 0.4268$ and
          $P({+}{-}) = P({-}{+}) = (1 - \cos 45^\circ)/4 \approx 0.0732$ --- the
          four sum to 1.
\end{itemize}

The standard summary is the \textbf{correlation function}
$E(\alpha,\beta) = P({+}{+}) + P({-}{-}) - P({+}{-}) - P({-}{+})$, which collapses
to a function of the angle \emph{difference} alone --- the state $\ket{\Phi^{+}}$ is
invariant when both analyzers rotate together:

\begin{equation}
    E(\alpha, \beta) = \cos 2\Delta, \qquad \Delta = \alpha - \beta.
    \label{eq:corr}
\end{equation}

\begin{table}[H]
\centering
\begin{tabular}{lcccc}
\toprule
analyzer difference $\Delta$ & $0^\circ$ & $22.5^\circ$ & $45^\circ$ & $90^\circ$ \\
\midrule
$E = \cos 2\Delta$ & $1$ & $\sqrt{2}/2 \approx 0.707$ & $0$ & $-1$ \\
\bottomrule
\end{tabular}
\caption{Correlation anchors, equation~\eqref{eq:corr}: perfect correlation,
the $\sqrt 2/2$ point, mutually unbiased analyzers, perfect anti-correlation.}
\label{tab:corr}
\end{table}

\begin{numberbox}
$E(0^\circ, 22.5^\circ) = \cos 45^\circ = \tfrac{\sqrt 2}{2} \approx \mathbf{0.707}$.
In \textbf{Preset D}, set $\alpha = 0^\circ$, $\beta = 22.5^\circ$, fire 10\,000
pairs, and read the measured $E$ against the theory overlay. Then set the dials equal
($\Delta = 0$) and confirm \textbf{100\% agreement} --- against the classical-wave
overlay's \textbf{75\%} --- while each side's own counter stays at 50/50. At
$\Delta = 45^\circ$, $E = 0$: the two analyzers have become mutually unbiased ---
Lab 6's key concept reappearing as zero correlation.
\end{numberbox}

Why is this \emph{deep} and not just pretty? Try the obvious classical story: each
pair is emitted with a shared random polarization $\lambda$, and each analyzer
responds by classical Malus, $\cos^2$ of the angle to $\lambda$. Average over
uniform $\lambda$: expanding each factor via $\cos^2 x = (1 + \cos 2x)/2$, the
single-cosine averages vanish, and the lone surviving average is
$\langle \cos 2(\alpha - \lambda)\cos 2(\beta - \lambda)\rangle_\lambda
= \cos(2\Delta)/2$, so

\begin{equation}
    P_{\text{cl}}({+}{+}) = \frac14 + \frac18\cos 2\Delta
    \qquad\Longrightarrow\qquad
    E_{\text{cl}}(\Delta) = \frac{\cos 2\Delta}{2}
    \label{eq:classical}
\end{equation}

--- the same shape, but \textbf{half the size}. At matched analyzers the classical
wave model manages only 75\% agreement ($E_{\text{cl}}(0) = \tfrac12$) against the
observed 100\%. This is not a philosophical complaint --- it is a
\textbf{device-verifiable factor of 2}, and the simulator's overlay draws both
curves so you can watch the measured points choose the quantum one. More
sophisticated classical models (predetermined answers for every angle) can patch the
matched-angle agreement but then fail elsewhere on the $\cos 2\Delta$ curve --- that
is \textbf{Bell's theorem} (Bell, 1964; Clauser--Horne--Shimony--Holt, 1969), which
we cite rather than derive: no local classical story reproduces the full curve. What
this lab establishes at the correlation level is the honest core: the measured curve
is $\cos 2\Delta$, and the classical wave prediction misses it by a factor of 2.

\begin{industrybox}
\textbf{Entanglement as infrastructure.} The pair source you are firing is the
physical basis of \textbf{entanglement-based QKD} (the BBM92 protocol): Alice and
Bob each receive one photon of a pair, measure in randomly switched bases, and the
matched-angle perfect correlations play the role that Alice's prepared states played
in Lab 6's BB84 --- with the added property that the correlations themselves certify
that no third party is entangled with the key. The same correlations, measured
across 1\,200 km by the \textbf{Micius} satellite, are the record distance for
entanglement distribution --- Part D at continental scale.
\end{industrybox}

\newpage

\section{Part E: Guided Activities}

\begin{simulatorbox}
All activities below run on the Optical Bench of the lab's interactive page
(\texttt{lab7\_optics.html}); the preset to use is named in each activity. All runs
are seeded --- the same seed always reproduces the same counts. Activity 6 uses the
companion Qiskit notebook:\\[4pt]
\url{https://colab.research.google.com/github/LMrilo/quantum-labs-2026/blob/main/lab7/lab7_optics.ipynb}
\end{simulatorbox}

\subsection{Activity 1: Indivisibility Statistics (Preset A)}

\begin{studentbox}
Fire 100, then 1\,000, then 10\,000 photons (same seed, resetting in between).
\begin{enumerate}
    \item How does the deviation of the detector-B fraction from $1/2$ shrink as $n$
          grows? Compare with the statistical scale $\sqrt{1/(4n)}$.
    \item Confirm the coincidence counter reads \emph{exactly} 0 in every run --- not
          just small. Why is this exact even though the simulation is random?
          (Equation~\eqref{eq:anticorr}.)
    \item What would a classical wave-pulse model predict for the coincidence
          counter?
\end{enumerate}
\end{studentbox}

\subsection{Activity 2: Mapping the Fringe (Preset B1)}

\begin{studentbox}
Use the $\varphi$ slider and the sweep button.
\begin{enumerate}
    \item Verify the three anchor points: $\varphi = 0^\circ \rightarrow$ 100/0,
          $\varphi = 60^\circ \rightarrow$ 75/25,
          $\varphi = 180^\circ \rightarrow$ 0/100 (Table~\ref{tab:mz}).
    \item Run the $\varphi$ sweep and check the measured dots trace
          $\cos^2(\varphi/2)$. Where on the curve does a small phase change produce
          the largest change in counts? (This is where an interferometric sensor
          operates.)
    \item Explain in one sentence why a photon that ``chose one arm at BS1'' cannot
          produce the $\varphi = 0$ result.
\end{enumerate}
\end{studentbox}

\subsection{Activity 3: Buying Back the Fringes (Preset B2)}

\begin{studentbox}
The complementarity dial: visibility $\mathcal{V} = s$.
\begin{enumerate}
    \item At $s = 0.5$, verify the fringe extremes 0.75 / 0.25 (fire at
          $\varphi = 0^\circ$ and $180^\circ$) and read the visibility card
          $\rightarrow$ 0.5.
    \item Run the $\varphi$ sweep at $s = 1$, $s = 0.5$, $s = 0$ and compare the
          three curves' amplitudes.
    \item Nothing ever reads the marker in this simulation. State precisely what
          kills the interference --- and connect it to why Eve in Lab 6 cannot
          listen without leaving a trace.
\end{enumerate}
\end{studentbox}

\subsection{Activity 4: Three Polarizers and the Staircase (Preset C)}

\begin{studentbox}
\begin{enumerate}
    \item Reproduce the pair of experiments: crossed $0^\circ/90^\circ \rightarrow$
          exactly 0; insert $45^\circ \rightarrow$ 12.5\%. Where do the absorbed
          photons go in each configuration? (Watch the per-polarizer absorption
          counts.)
    \item With the $\ket{H}$ source and a single polarizer at $60^\circ$, verify
          Malus: 25\% pass.
    \item Set the staircase $0^\circ/30^\circ/60^\circ/90^\circ$ and check
          $\tfrac12\,(3/4)^3 = 27/128 \approx 21.1\%$. Explain why \emph{more}
          absorbing filters transmit \emph{more} light, using the phrase
          ``post-measurement state''.
\end{enumerate}
\end{studentbox}

\subsection{Activity 5: The Factor of 2 (Preset D --- Optional)}

\begin{studentbox}
\begin{enumerate}
    \item Set $\alpha = \beta = 20^\circ$: verify each side's own counter is
          $\approx$ 50/50 while agreement is 100\%. Change the common angle; does
          anything change?
    \item Set $\alpha = 0^\circ$, $\beta = 22.5^\circ$ and verify
          $E \approx 0.707$ against the overlay. At $\Delta = 45^\circ$, why is
          $E = 0$? (Lab 6 vocabulary: what are the two analyzers to each other?)
    \item Run the $\Delta$ sweep. The dashed curve is the best classical-wave story.
          Identify the angle where the quantum/classical gap is largest, and say
          what an experimenter would conclude from data landing on the solid curve.
\end{enumerate}
\end{studentbox}

\subsection{Activity 6: The Optics in Qiskit}

The companion notebook rebuilds this guide's optics as quantum circuits: the beam
splitter as a one-qubit gate on the dual-rail path qubit, the Mach--Zehnder as a
three-gate circuit swept over $\varphi$ against $\cos^2(\varphi/2)$, which-path
marking as a controlled-NOT onto an ancilla with the visibility sweep
$\mathcal{V} = s$, the polarization experiments with Procedure 1 and Procedure 2
side by side, and (optional) the entangled-pair correlation curve against
$\cos 2\Delta$. Its core circuit is three gates long:

\begin{lstlisting}[language=Python]
def mach_zehnder(phi: float) -> QuantumCircuit:
    """Mach-Zehnder on the dual-rail path qubit: BS - phase - BS."""
    qc = QuantumCircuit(1)
    qc.rx(-np.pi / 2, 0)   # beam splitter  B  (the i-on-reflection 50/50)
    qc.p(phi, 0)           # phase shifter  Phi(phi)  on arm B
    qc.rx(-np.pi / 2, 0)   # beam splitter  B
    qc.measure_all()
    return qc
\end{lstlisting}

\noindent D1 collects the photons observed in path B (outcome 1), so the counts of
outcome 1 trace $P(D1) = \cos^2(\varphi/2)$ across the sweep. All notebook
experiments are seeded and reproducible; expected statistics and tolerances are
stated inline. A compact optional section runs the interferometer on real IBM
hardware.

\begin{warningbox}
\textbf{Qiskit bit-ordering note.} Qiskit orders multi-qubit bitstrings
little-endian (rightmost bit = qubit $q_0$), while this lab series writes states
big-endian (qubit 0 leftmost). Unlike Lab 6, where every register held a single
qubit, this lab genuinely needs the conversion: two-qubit registers appear in the
which-path model (the $\ket{\text{path}\;\text{marker}}$ states) and in Part D (the
photon pairs), so \emph{every} displayed bitstring must be converted to the
leftmost-first ordering. The notebook shows exactly how, and applies it to every
histogram.
\end{warningbox}

\subsection{Activity 7: Quick Check}

\begin{reflectionbox}
\textbf{Question 1:} A single photon in state $(\ket{A} + i\ket{B})/\sqrt 2$ meets
two detectors, one per path. What is the probability both click?
\begin{itemize}
    \item[(a)] $1/4$ \qquad (b) exactly $0$ \qquad (c) small but nonzero
\end{itemize}

\textbf{Question 2:} A Mach--Zehnder interferometer at $\varphi = 0$ is built from
two 50/50 beam splitters. What fraction of photons reaches D1?
\begin{itemize}
    \item[(a)] 50\% --- each splitter is a fair coin.
    \item[(b)] 100\% --- the two path amplitudes interfere.
    \item[(c)] 75\%.
\end{itemize}

\textbf{Question 3:} A photon passes a polarizer at $45^\circ$ (Procedure 1) and, in
a second experiment, a half-waveplate at $22.5^\circ$ + polarizing beam splitter
(Procedure 2), in both cases with the ``+'' outcome. The outcome
\emph{probabilities} were identical. The exiting photon is\dots
\begin{itemize}
    \item[(a)] the same state in both cases --- probabilities fix the state.
    \item[(b)] $\ket{+}$ after Procedure 1, $\ket{H}$ after Procedure 2 --- overlap
               only $1/2$.
    \item[(c)] unpolarized in both cases.
\end{itemize}
\end{reflectionbox}

\section{Resources}

\begin{itemize}
    \item Grangier, P., Roger, G. \& Aspect, A., ``Experimental evidence for a
          photon anticorrelation effect on a beam splitter: a new light on
          single-photon interferences,'' \textit{Europhys.\ Lett.}\ \textbf{1}, 173
          (1986). --- Part A's anti-correlation and Part B's single-photon
          interference in one experiment.
    \item Zeilinger, A., ``General properties of lossless beam splitters in
          interferometry,'' \textit{Am.\ J.\ Phys.}\ \textbf{49}, 882 (1981). ---
          Beam-splitter unitarity and the phase convention of Part A.
    \item Loudon, R., \textit{The Quantum Theory of Light}, 3rd ed., Oxford
          University Press, 2000, \S\S 5--6. --- Beam splitters, single-photon
          states, and the $g^{(2)}$ correlation cited in Part A.
    \item Englert, B.-G., ``Fringe visibility and which-way information: an
          inequality,'' \textit{Phys.\ Rev.\ Lett.}\ \textbf{77}, 2154 (1996). ---
          The duality relation $\mathcal{V}^2 + \mathcal{D}^2 \leq 1$ behind
          Part B's remark.
    \item Kwiat, P. G. et al., ``New high-intensity source of
          polarization-entangled photon pairs,'' \textit{Phys.\ Rev.\ Lett.}\
          \textbf{75}, 4337 (1995). --- SPDC pair sources, Part D's physical
          platform.
    \item Bell, J. S., ``On the Einstein Podolsky Rosen paradox,''
          \textit{Physics} \textbf{1}, 195 (1964); Clauser, J. F., Horne, M. A.,
          Shimony, A. \& Holt, R. A., \textit{Phys.\ Rev.\ Lett.}\ \textbf{23}, 880
          (1969). --- Cited in Part D to back the ``no classical story'' claim;
          derivations beyond this lab's scope.
    \item Nielsen, M. A. \& Chuang, I. L., \textit{Quantum Computation and Quantum
          Information}, Cambridge University Press, 2010, \S 7.4. --- Optical
          photon quantum computation (the dual-rail encoding of Part A).
    \item Kok, P. et al., ``Linear optical quantum computing with photonic
          qubits,'' \textit{Rev.\ Mod.\ Phys.}\ \textbf{79}, 135 (2007). --- The
          LOQC/KLM context for Part A's industry anchor.
    \item Hong, C. K., Ou, Z. Y. \& Mandel, L., \textit{Phys.\ Rev.\ Lett.}\
          \textbf{59}, 2044 (1987). --- Hong--Ou--Mandel interference: further
          reading, beyond this lab's scope.
    \item Lab 5: \textit{Quantum Measurements} and Lab 6: \textit{Quantum Key
          Distribution (BB84)} (this series). --- Born rule, Procedures 1 and 2,
          mutual unbiasedness, and the measurement--disturbance--security bridge.
    \item Companion Jupyter notebook:
          \texttt{quantum-labs-2026/lab7/lab7\_optics.ipynb} (Colab link in
          Part E).
\end{itemize}

\end{document}
